\subsection{De karakteristieke veelterm van een lineaire operator}

\begin{lemma}[6.1.1]
    Zij $V$ een eindig-dimensionale $K$-vectorruimte en $f \in \operatorname{End}(V)$. Beschouw twee matrixvoorstellingen $A$ en $B$ van $f$ ten opzichte van verschillende basissen van $V$. Dan is $\det A = \det B$.
\end{lemma}
\begin{proof}
    Wegens (Stelling 4.2.5) bestaat er een inverteerbare matrix $P \in \operatorname{GL}_n(K)$ zodat:
    \begin{align*}
        A &= P^{-1}BP\\
        \det(A) &= \det(P^{-1}BP) = \det(P^{-1})\det(B)\det(P)\\
        &= \det(B)\det(P^{-1})\det(P) = \det(B)\det(I)\\
        &= \det(B)
    \end{align*}
\end{proof}

\begin{lemma}[6.1.4]
    Zij $V$ een eindig-dimensionale vectorruimte over een veld $K$, en zij $f \in \operatorname{End}_K(V)$ een lineaire operator, met matrixvoorstelling $A$ (ten opzichte van een willekeurige basis voor $V$). Dan is
    \[ \chi_f(x) := \det(xI_n - A) \in K[x] \]
    een monische veelterm van graad $n$ in de variabele $x$, en deze veelterm hangt niet af van de keuze van de basis voor $V$.
\end{lemma}
\begin{proof}
    We tonen eerst aan dat \(\chi_f(x)\) basis onafhankelijk is.
    Wegens (Stelling 4.2.5)  $\exists P \in \operatorname{GL}_n(K)$ zodat:
    \[A = P^{-1}BP\]
    \begin{align*}
        \det(xI_n - A) 
        &= \det(xI_n - P^{-1}BP)\\
         &= \det(P(xI_n - P^{-1}BP)P^{-1})\\
         &= \det(xPI_nP^{-1} - PP^{-1}BPP^{-1})\\
         &= \det((xI_n - B)) \\
    \end{align*}
    Waarbij we gebruik hebben gemaakt van (stelling 4.2.5) en de eigenschap dat \(I_n\) commuteerd.

    Doordat in de determinant alle factoren met \(x\) op de diagonaalelementen staan,
    is de term met de hoogste graad \(\prod_{i = 1}^n (x - a_{ii})\).

    Hier is de coëfficiënt van de term met de hoogste graad duidelijk 1.
\end{proof}

\begin{lemma}[6.1.6]
    Zij $A = (a_{ij}) \in M_n(K)$ met karakteristieke veelterm
    \[ \chi_A(x) = x^n + c_{n-1}x^{n-1} + c_{n-2}x^{n-2} + \dots + c_1x + c_0. \]
    Dan is $c_0 = (-1)^n \det A$ en $c_{n-1} = -\sum_{i=1}^n a_{ii}$.
\end{lemma}
\begin{proof}
    Stel \(x = 0\), dan is \(\chi_A(x=0) = c_0\).

    \(\implies c_0 = \det(-A) = (-1)^n \det(A)\)

    Waarbij we gebruik hebben gemaakt van (lemma 5.4.10 (iii)).



    $$\det(xI_n - A) = \sum_{\sigma \in S_n} \operatorname{sgn}(\sigma) \prod_{i=1}^n B_{i, \sigma(i)} $$ 

    De hoogste graad is \(n\) (lemma 6.1.4), deze komt overeen met de permutatie \(\sigma = id\).
    Doordat een permutatie 2 elementen wisselt, is de hoogste graad wanneer \(\sigma \neq id, \, n-2\).

    \(\implies\) de \(n-1\) term moet van de diagonaal-termen komen (\(\sigma = id\)).

    \[\prod_{i=1}^n (x - a_{ii}) = x^n - \left(\sum_{i=1}^n a_{ii}\right)x^{n-1} + \dots\]

    De enige \(n-1\) term die we hieruit kunnen halen is \(-\sum_{i=1}^{n} a_{ii}\).


\end{proof}

\subsection{Eigenwaarden en eigenvectoren}

\begin{stelling}[6.2.5]
    \begin{enumerate}
        \item[(i)] Zij $A \in M_n(K)$. Een scalair $\lambda \in K$ is een eigenwaarde van $A$ als en slechts als $\chi_A(\lambda) = \det(\lambda I_n - A) = 0$.
        \item[(ii)] Zij $V$ een $n$-dimensionale $K$-vectorruimte en $f \colon V \to V$ een lineaire operator. Een scalair $\lambda \in K$ is een eigenwaarde van $f$ als en slechts als $\chi_f(\lambda) = \det(\lambda \mathbf{1} - f) = 0$.
    \end{enumerate}
\end{stelling}
\begin{proof}(i)
    \([\Rightarrow]\) Als \(\lambda\) een eigenwaarde van \(A\), dan geldt:
    \(f(v) = Av =\lambda v, \, 0 \neq v \in V\).

    \begin{align*}
        Av &= \lambda v\\
        \implies 0 &= (\lambda I_n - A)v\\
        \implies v &\in \ker(\lambda I_n - A)\\
        \implies &\dim(\ker(\lambda I_n - A)) > 0\\
        \implies &rk((\lambda I_n - A)) < n \\
        \implies &\det(\lambda I_n - A) = 0
    \end{align*}

    \([\Leftarrow]\) omgekeerd als \(\det(\lambda I_n - A) = 0\):
    \begin{align*}
        \implies &rk((\lambda I_n - A)) < n \\
        \implies &\dim(\ker(\lambda I_n - A)) > 0\\
        \implies & \exists v (\neq 0) \in V: (\lambda I_n - A)v = 0 ,\\
        \implies & \exists v (\neq 0) \in V:  Av = \lambda v
    \end{align*}
\end{proof}
\begin{proof}(ii)
    Kies een willekeurige basis \(\mathcal{B}\) van \(V\) en \(\beta \colon V \to K^n\) het  coördinatenisomorfisme. 
    Als \(A := A_{f,\mathcal{B},\mathcal{B}}\) de matrixvoorstelling van \(f\) tov \(\mathcal{B}\).

    Wegens Opmerking 4.1.5 (ii) geldt \( \forall v \in V: \beta(f(v)) = A\beta(v)\).
    \( \forall v \in V \setminus \{0\}\) en \(\forall \lambda \in K\) gelden:
    \[ f(v) = \lambda v \iff \beta(f(v)) = \beta(\lambda v) \iff A\beta(v) = \lambda \beta(v). \]

    \(\implies v \in V \setminus \{0\}\) een eigenvector is van \(f\) met eigenwaarde \(\lambda \iff \beta(v) \in K^n \setminus \{0\}\) een eigenvector is van \(A\) met eigenwaarde \(\lambda\).

    Bijgevolg is \(\lambda\) een eigenwaarde van \(f \iff \lambda\) een eigenwaarde is van \(A\).
    Wegens (i) geldt:
    \[ \chi_f(\lambda) = \det(\lambda \mathbf{1} - f) = \det(\lambda I_n - A) = 0. \]
\end{proof}

\begin{gevolg}[6.2.6]
    Een lineaire operator $f$ op een $n$-dimensionale $K$-vectorruimte $V$ (of een matrix $A \in M_n(K)$) heeft hoogstens $n$ eigenwaarden.
\end{gevolg}
\begin{proof}
    Uit (lemma 6.1.4) volgt dat \(\chi_f(x)\) een monische veelterm is van graad \(n\).

    Uit (stelling 6.2.5) (ii) is \(\lambda\) een eigenwaarde  
    \(\iff \chi_f(\lambda) = 0. \)

    Een veelterm in \(K\) heeft hoogstens \(n\) wortels, waaruit het gevraagde geldt. 

\end{proof}

\subsection{Diagonaliseren van operatoren}

\begin{stelling}[6.3.1]
    Een lineaire operator $f$ op een eindig-dimensionale $K$-vectorruimte heeft een diagonaalmatrix als matrixvoorstelling als en slechts als $V$ een basis heeft bestaande uit eigenvectoren voor $f$.
\end{stelling}
\begin{proof}
    $[\Rightarrow]$ Stel dat $f$ tov de basis $\mathcal{B} = \{b_1, \dots, b_n\}$ een diagonaalmatrix $A = \operatorname{diag}(\lambda_1, \dots, \lambda_n)$ als voorstelling heeft.
     De $i$-de kolom van $A$ is de coördinatenvector van het beeld $f(b_i)$ tov de basis $\mathcal{B}$.
    $ \forall 1 \le i \le n$:
    \[ f(b_i) = 0b_1 + \dots + 0b_{i-1} + \lambda_i b_i + 0b_{i+1} + \dots + 0b_n = \lambda_i b_i. \]
    Aangezien $\mathcal{B}$ een basis is, geldt dat $b_i \neq 0$ voor alle $i$.
    Dit betekent per definitie dat elke $b_i$ een eigenvector is van $f$ met bijhorende eigenwaarde $\lambda_i$.
    Bijgevolg is $\mathcal{B}$ een basis bestaande uit eigenvectoren van $f$.

    $[\Leftarrow]$ Zij omgekeerd $\mathcal{B} = \{b_1, \dots, b_n\}$ een basis van $V$ bestaande uit eigenvectoren voor de operator $f \in \operatorname{End}(V)$.
    $\exists \forall 1 \le i \le n$ een scalair $\lambda_i \in K$ waarvoor:
    \[ f(b_i) = \lambda_i b_i. \]
    tov de basis $\mathcal{B}$:
    \(f(b_i) = 0b_1 + \dots + \lambda_i b_i + \dots + 0b_n. \)
    Wanneer we deze coördinatenvectoren als kolommen in de matrixvoorstelling van $f$ plaatsen, bekomen we de diagonaalmatrix $\operatorname{diag}(\lambda_1, \dots, \lambda_n)$.
\end{proof}

\begin{lemma}[6.3.3]
    Een matrix $A \in M_n(K)$ is diagonaliseerbaar als en slechts als er een matrix $P \in \operatorname{GL}_n(K)$ bestaat zodat $P^{-1}AP$ een diagonaalmatrix is.
\end{lemma}
\begin{proof}
    \([\Rightarrow]\) \(A\) is diagonaliseerbaar:
    
    Er bestaat een basis \(\mathcal{B} = (b_1, \dots, b_n)\) uit eigenvectoren.

    Als \(P = (b_1, \dots, b_n)\), dan is \(AP = (\lambda_1 b_1, \dots, \lambda_n b_n) = \operatorname{diag}(b_1, \dots, b_n) \cdot (\lambda_1,\dots,\lambda_n)\). 
    En \(P^{-1}AP = \operatorname{diag}(\lambda_1,\dots\lambda_n)\).

    \([\Leftarrow]\) Omgekeerd, \(P^{-1}AP= \operatorname{diag}(\lambda_1,\dots\lambda_n)\) :

    Beschouw basis \(\mathcal{B} = (b_1, \dots, b_n)\) met \(b_i := Pe_i\)

    Dan \(P^{-1}A(b_1, \dots , b_n)= \operatorname{diag}(\lambda_1,\dots\lambda_n)\)

    \(\implies A(b_1, \dots , b_n)= P \operatorname{diag}(\lambda_1,\dots\lambda_n)\)

    Het is duidelijk dat \(L_A \) tov \(\mathcal{B}\) de diagonaalmatrix \(P^{-1}AP\) is.
    
\end{proof}

\begin{stelling}[6.3.4]
    Zij $V$ een $n$-dimensionale vectorruimte, en zij $f \in \operatorname{End}(V)$.
    \begin{enumerate}
        \item[(i)] De eigenvectoren van $f$ die horen bij verschillende eigenwaarden, zijn lineair onafhankelijk van elkaar.
        \item[(ii)] Als $f$ precies $n$ verschillende eigenwaarden heeft, dan is $f$ diagonaliseerbaar.
        \item[(iii)] Als $f$ diagonaliseerbaar is, dan is $V$ de directe som van de eigenruimten van $f$ horende bij de verschillende eigenwaarden van $f$.
    \end{enumerate}
\end{stelling}
\begin{proof}(i)
    Stel dat dit niet zo was. \(\mu_1v_1 + \dots + \mu_n v_m = 0\) met \(m\) minimaal zodat alle termen niet nul zijn.


    \(f(\mu_1v_1 + \dots + \mu_m v_m ) = \lambda_1 \mu_1 v_1 + \dots + \lambda_m \mu_m v_m = 0\)

    Aangezien \(\mu_1v_1 + \dots + \mu_m v_m = 0\), dan is ook \(\lambda_1\mu_1v_1 + \dots + \lambda_1\mu_m v_m = 0\) een lineaire combinatie.

    Nu is \( \lambda_1 \mu_1 v_1 + \dots + \lambda_m \mu_m v_m - (\lambda_1\mu_1v_1 + \dots + \lambda_1\mu_m v_m) = 0 \)

    \(\implies (\lambda_2-\lambda_1)\mu_2 v_2 + \dots + (\lambda_n-\lambda_1)\mu_m v_m = 0 \).

    Maar deze heeft minder termen dan ze minimale \(m\) die we gesteld hadden.

    Dus de eigenvectoren zijn lineair onafhankelijk.

\end{proof}
\begin{proof}(ii)
    Uit (i) zijn de eigenvectoren lineair onafhankelijk. Uit (stelling 2.4.12 (i)) vormen de eigenvectoren een basis voor V.
    Uit (stelling 6.3.1) is \(f\) diagonaliseerbaar. 
\end{proof}
\begin{proof}(iii)
    Doordat de eigenruimten die bij verschillende eigenvectoren horen lineair onafhankelijk zijn,
    is de van gelijk welke eigenruimte van verschillende eigenvectoren \(\{0\}\).
\end{proof}
\begin{proof}(iii)
    Zij $\lambda_1, \dots, \lambda_\ell$ de verschillende eigenwaarden van $f$.
    Aangezien $f$ diagonaliseerbaar is, bestaat er wegens Stelling 6.3.1 een basis $\mathcal{B}$ voor $V$ die volledig bestaat uit eigenvectoren van $f$. 
    
    We kunnen deze basis $\mathcal{B}$ schrijven als de disjuncte unie van deelverzamelingen:
    \(\mathcal{B} = \mathcal{B}_1 \cup \dots \cup \mathcal{B}_\ell \)
    waarbij elke $\mathcal{B}_i$ precies de basisvectoren uit $\mathcal{B}$ bevat die horen bij de eigenwaarde $\lambda_i$.
    
    Per definitie van de eigenruimten geldt dat de eigenruimte $V_i$ behorende bij $\lambda_i$ gelijk is aan:
    \(V_i = \operatorname{span}(\mathcal{B}_i). \)
    
    Omdat $\mathcal{B}$ een basis is voor $V$ en de disjuncte unie is van de verzamelingen $\mathcal{B}_i$, volgt uit de transformatiestelling voor basissen (Stelling 2.5.5, recursief toegepast op $\ell$ delen) onmiddellijk dat:
    \[ V = \operatorname{span}(\mathcal{B}_1) \oplus \dots \oplus \operatorname{span}(\mathcal{B}_\ell) \]
    \[ \implies V = V_1 \oplus \dots \oplus V_\ell. \]
\end{proof}

\begin{lemma}[6.3.9]
    Zij $V$ een $K$-vectorruimte en $f \colon V \to V$ een lineaire operator met karakteristieke veelterm
    \[ \chi_f(x) = \prod_{i=1}^t (x - \lambda_i)^{n_i} \quad \text{met } \lambda_1, \dots, \lambda_t \in K. \]
    Voor iedere $i$ is $\dim(\operatorname{ker}(\lambda_i \mathbf{1}_V - f)) \le n_i$. Anders gezegd, de meetkundige multipliciteit van een eigenwaarde is steeds kleiner dan of gelijk aan de algebraïsche multipliciteit van deze eigenwaarde.
\end{lemma}
\begin{proof}
    We noteren de eigenruimte behorende bij de eigenwaarde $\lambda_i$ met $V_i = \ker(\lambda_i \mathbf{1}_V - f)$.
     Zij $r_i := \dim(V_i)$ de meetkundige dimensie van $\lambda_i$.
      We tonen aan dat $(x-\lambda_i)^{r_i}$ een deler is van $\chi_f(x)$,
      wat direct impliceert dat $r_i \le n_i$.

    Stel $r = r_i$ en $\lambda = \lambda_i$.
     Kies een basis $\{v_1, \dots, v_r\}$ voor $V_i$.
      Wegens Gevolg 2.4.8(i) breiden we deze lineair onafhankelijke verzameling uit tot een basis voor de totale ruimte $V$:
    \[ \mathcal{B} = \{v_1, \dots, v_r, w_1, \dots, w_s\} \quad \text{met } s = n - r. \]
    Omdat $f(v_j) = \lambda v_j , \, \forall 1 \le j \le r$, 
    ziet de matrixvoorstelling $A$ van $f$ tov de basis $\mathcal{B}$ er als volgt uit:
    \[ A = \begin{pmatrix} \lambda I_r & B \\ 0 & C \end{pmatrix} \]
    voor zekere matrices $B \in M_{r,n-r}(K)$ and $C \in M_{n-r}(K)$. Er geldt nu dat:
    \[ \chi_f(x) = \det(x I_n - A) = \det \begin{pmatrix} (x - \lambda) I_r & -B \\ 0 & x I_{n-r} - C \end{pmatrix}. \]
    Wanneer we deze determinant $r$ opeenvolgende malen ontwikkelen naar de eerste kolom, bekomen we:
    \[ \chi_f(x) = (x - \lambda)^r \det(x I_{n-r} - C). \]
    Hieruit volgt dat $(x - \lambda)^r$ een deler is van $\chi_f(x)$, waardoor de meetkundige dimensie $r$ hoogstens gelijk kan zijn aan de algebraïsche multipliciteit $n_i$.
\end{proof}

\begin{stelling}[6.3.10]
    Zij $V$ een $K$-vectorruimte en $f \colon V \to V$ een lineaire operator met karakteristieke veelterm
    \[ \chi_f(x) = \prod_{i=1}^t (x - \lambda_i)^{n_i} \quad \text{met } \lambda_1, \dots, \lambda_t \in K. \]
    Zij $V_i = \operatorname{ker}(\lambda_i \mathbf{1}_V - f)$ de eigenruimte behorende bij de eigenwaarde $\lambda_i$, voor elke $i$. Dan is $f$ diagonaliseerbaar als en slechts als voor elke eigenwaarde $\lambda_i$ geldt dat $\dim(V_i) = n_i$, of met andere woorden, als en slechts als voor iedere eigenwaarde de algebraïsche multipliciteit gelijk is aan de meetkundige multipliciteit.
\end{stelling}
\begin{proof}
    $\operatorname{deg}(\chi_f(x))=n \implies n = \dim V = n_1 + \dots + n_t$.
    Voor elke eigenruimte $V_i$ kiezen we een basis $\mathcal{B}_i$.
    Wegens Stelling 6.3.4(i) is $\mathcal{C} = \mathcal{B}_1 \cup \dots \cup \mathcal{B}_t$ lineair onafhankelijk.

    \([\Leftarrow]\)
    Stel dat $\dim(V_i) = n_i, \, \forall i$.
     Aangezien de basissen $\mathcal{B}_i$ onderling disjunct zijn (ze horen bij verschillende eigenwaarden),
     is het totale aantal elementen in de verzameling $\mathcal{C}$ gelijk aan:
    \[ |\mathcal{C}| = \sum_{i=1}^t \dim(V_i) = \sum_{i=1}^t n_i = n . \]
    Omdat $\mathcal{C}$ een lineair onafhankelijke verzameling is met $n$ elementen in een $n$-dimensionale ruimte,
    vormt $\mathcal{C}$ een basis van $V$ bestaande uit eigenvectoren van $f$.
    Wegens Stelling 6.3.1 is $f$ bijgevolg diagonaliseerbaar.

    \([\Rightarrow]\) 
    Veronderstel dat $\dim(V_i) < n_i$ voor een bepaalde index $i$.
     Aangezien voor elke andere index wegens Lemma 6.3.9 geldt dat $\dim(V_j) \le n_j$,
    is het totale aantal elementen in de lineair onafhankelijke verzameling $\mathcal{C}$ strikt kleiner dan $n$:
    \[ |\mathcal{C}| = \sum_{j=1}^t \dim(V_j) < \sum_{j=1}^t n_j = n. \]
    Omdat elke eigenvector van $f$ in de span van $\mathcal{C}$ moet liggen,
    kunnen we nooit $n$ lineair onafhankelijke eigenvectoren in $V$ vinden.
    Er bestaat dus geen basis van eigenvectoren voor $V$, wat wegens Stelling 6.3.1 impliceert dat $f$ niet diagonaliseerbaar is.
\end{proof}